% 136-34(136쪽) — 정사각형·정오각형·정육각형에 대각선을 모두 그은 그림과 그 아래 $a_4=2$, $a_5=5$, $a_6=9$ 라벨
% (원문 「다음 그림」). 스캔 화질이 낮아 TikZ 로 다시 그림.
% 라벨은 원문 것만 둔다 — $a_4=2$ · $a_5=5$ · $a_6=9$ 와 오른쪽의 $\cdots$ 두 개.
% 원문처럼 육각형까지만 그려 구하는 $f(20)$(답)이 그림에서 읽히지 않게 한다.
\begin{tikzpicture}[line width=0.45pt, x=1cm, y=1cm, font=\small]
  % 정사각형 — 대각선 2개
  \coordinate (q1) at (0.00,0.05);
  \coordinate (q2) at (1.60,0.05);
  \coordinate (q3) at (1.60,1.65);
  \coordinate (q4) at (0.00,1.65);
  \draw[line width=0.7pt] (q1) -- (q2) -- (q3) -- (q4) -- cycle;
  \draw (q1) -- (q3);
  \draw (q2) -- (q4);
  \node at (0.80,-0.42) {$a_4=2$};
  % 정오각형 — 대각선 5개
  \foreach \i/\a in {1/90, 2/162, 3/234, 4/306, 5/18}{
    \coordinate (p\i) at ($(3.05,0.85)+(\a:0.95)$);
  }
  \draw[line width=0.7pt] (p1) -- (p2) -- (p3) -- (p4) -- (p5) -- cycle;
  \foreach \u/\v in {1/3, 1/4, 2/4, 2/5, 3/5}{
    \draw (p\u) -- (p\v);
  }
  \node at (3.05,-0.42) {$a_5=5$};
  % 정육각형 — 대각선 9개
  \foreach \i/\a in {1/90, 2/150, 3/210, 4/270, 5/330, 6/30}{
    \coordinate (h\i) at ($(5.45,0.85)+(\a:0.92)$);
  }
  \draw[line width=0.7pt] (h1) -- (h2) -- (h3) -- (h4) -- (h5) -- (h6) -- cycle;
  \foreach \u/\v in {1/3, 1/4, 1/5, 2/4, 2/5, 2/6, 3/5, 3/6, 4/6}{
    \draw (h\u) -- (h\v);
  }
  \node at (5.45,-0.42) {$a_6=9$};
  % 원문의 말줄임
  \node at (6.75,0.95) {$\cdots$};
  \node at (6.75,-0.42) {$\cdots$};
\end{tikzpicture}
